\documentclass[12pt]{article}
\usepackage[utf8]{inputenc}
\usepackage[T1]{fontenc}
\usepackage[margin=1in]{geometry}
\usepackage{setspace}
\usepackage{xcolor}
\usepackage[colorlinks=true, urlcolor=blue!60!black, linkcolor=blue!60!black]{hyperref}
\setlength{\parindent}{2em}
\setlength{\parskip}{1.25em}
\renewcommand{\baselinestretch}{1.0}
\begin{document}
\begin{center}
{\LARGE \textit{Geduldig v. Aiello}, 417 U.S. 484 (1974)\par}
\vspace{0.8em}
{\large Case Brief --- Four-Page Edition\par}
\end{center}
\vspace{1em}

\section{Memory Jogger}

Excluding normal pregnancy from California's employee disability insurance program was not sex discrimination under the Equal Protection Clause, because the program divided recipients into ``pregnant women and nonpregnant persons,'' triggering only rational basis review (\href{https://supreme.justia.com/cases/federal/us/417/484/}{Justia}).

\section{Detailed Case Facts}

Since 1946 California had run a mandatory, entirely self-supporting disability insurance system for private employees unable to work because of illness or injury not covered by workers' compensation. Employees contributed 1\% of wages up to an annual maximum, with no general tax subsidy, and received weekly benefits for up to 26 weeks. The program excluded disabilities lasting fewer than eight days without hospitalization, disabilities beyond 26 weeks, and certain court commitments. Section 2626 also excluded ``any injury or illness caused by or arising in connection with pregnancy up to the termination of such pregnancy and for a period of 28 days thereafter.''

Four working women disabled by pregnancy were denied benefits and sued. Three had abnormal complications (ectopic and tubal pregnancies and a miscarriage); Jacqueline Jaramillo had a normal pregnancy. After a state appellate court construed \S~2626 in \textit{Rentzer} to exclude only normal pregnancy, California paid the three abnormal-pregnancy claims, leaving only Jaramillo's claim live.

\section{Procedural History}

The consolidated suits were heard by a three-judge district court, which in \textit{Aiello v. Hansen} (N.D. Cal. 1973) granted summary judgment for plaintiffs, holding that the exclusion lacked a rational and substantial relationship to a legitimate state purpose. On direct appeal, the Supreme Court stayed that judgment, heard argument on March 26, 1974, and reversed on June 17, 1974.

\section{Judicial Votes}

Reversed, \textbf{6--3}.
\begin{itemize}
\item \textbf{Majority:} Justice Stewart, joined by Chief Justice Burger and Justices White, Blackmun, Powell, and Rehnquist.
\item \textbf{Concurrences:} None.
\item \textbf{Dissent:} Justice Brennan, joined by Justices Douglas and Marshall.
\end{itemize}

\section{Holding}

The abnormal-pregnancy claims were moot. Excluding normal pregnancy from a comprehensive, self-supporting disability insurance program is not a sex-based classification triggering heightened scrutiny; absent a showing that a pregnancy distinction is a mere pretext for invidious sex discrimination, it receives rational basis review and is valid if it rationally furthers legitimate fiscal and administrative interests. Footnote 20's conceptual rationale became the decision's most consequential rule, later cited in \textit{General Electric Co. v. Gilbert}, \textit{Dobbs}, and \textit{United States v. Skrmetti}.

\section{Analysis of Opinions}

\textbf{Majority (Stewart).} The program was a social-welfare insurance scheme subject to deferential review under \textit{Williamson v. Lee Optical} and \textit{Dandridge v. Williams}. The fund paid out 90\% to 103\% of annual revenues, and covering normal pregnancy would cost an estimated \$120 million to \$131 million a year (a 33\% to 36\% increase), forcing the State to abandon self-support, cut benefits, or raise contributions in a way that would burden low-income workers. A State need not ``choose between attacking every aspect of a problem or not attacking the problem at all''; excluding pregnancy was no different in principle from excluding short or very long disabilities. To avoid heightened scrutiny under \textit{Reed} and \textit{Frontiero}, Stewart reasoned that the program insured the same risks for both sexes, and in Footnote 20 he held that pregnancy classifications are not sex classifications absent proof of pretext; women contributed about 28\% of the fund but received 38\% of benefits (n. 21).

\textbf{Dissent (Brennan).} The statute created a gender-based double standard warranting strict scrutiny. The fund covered virtually every other disabling condition regardless of cost, voluntariness, or ``normalcy''---heart attacks, cosmetic surgery, and sex-linked conditions such as prostatectomies and gout---yet singled out a disability unique to women. Citing the EEOC's Title VII guidelines, he argued that saving money cannot justify an invidious classification, especially when sex-neutral adjustments to the contribution rate could preserve solvency.

\section{Examples: Future Applications}

\begin{enumerate}
\item \textbf{Same-Side (Upheld): Abortion Funding.} A Medicaid program that funds childbirth but not elective abortion is not a sex classification absent pretext, under \textit{Geduldig} and \textit{Dobbs}.
\item \textbf{Same-Side (Upheld): Treatment Ban.} Extending \textit{Geduldig}'s ``lack of identity'' logic, \textit{Skrmetti} held that a ban on puberty blockers for gender dysphoria, but not other diagnoses, gets only rational basis review.
\item \textbf{Opposite-Side (Invalid): Forced Leave at Five Months.} A mandatory leave policy resting on stereotypes about pregnant women's capacity to work is pretextual sex discrimination and violates Title VII as amended by the Pregnancy Discrimination Act.
\item \textbf{Opposite-Side (Struck Down): Refusing Previously Pregnant Applicants.} A university that will not hire women who have been pregnant, while hiring fathers, uses pregnancy as a transparent pretext and fails intermediate scrutiny.
\item \textbf{Fence-Sitter: Contraceptive Exclusion.} A public health plan excludes contraceptives but covers erectile-dysfunction drugs. The State says it classifies by condition to contain costs; plaintiffs invoke Footnote 20's pretext exception.
\end{enumerate}

\section{Critique}

\textbf{Defenders and formalist doubts.} Originalist and formalist scholars defend the decision's refusal to expand heightened scrutiny beyond facial classifications, leaving fiscal trade-offs in benefit programs to legislatures. Yet even some formalists call the ``pregnant women and nonpregnant persons'' syllogism strained, since a trait exclusive to one sex is a perfect proxy for that sex.

\textbf{Feminist and progressive critics.} \textit{Geduldig} is among the most criticized equal protection decisions. Scholars such as Reva Siegel and Cary Franklin argue it enshrined an ``invisible male norm,'' treating male-typical disabilities as neutral and pregnancy as an anomalous extra cost. Congress repudiated its Title VII sequel, \textit{Gilbert}, in the Pregnancy Discrimination Act of 1978, and critics see its revival in \textit{Dobbs} and \textit{Skrmetti} as a setback for sex equality.

\section{Key Quotations}

\begin{quote}
``While it is true that only women can become pregnant, it does not follow that every legislative classification concerning pregnancy is a sex-based classification like those considered in \textit{Reed}, \textit{supra}, and \textit{Frontiero}, \textit{supra}. Normal pregnancy is an objectively identifiable physical condition with unique characteristics. Absent a showing that distinctions involving pregnancy are mere pretexts designed to effect an invidious discrimination against the members of one sex or the other, lawmakers are constitutionally free to include or exclude pregnancy from the coverage of legislation such as this on any reasonable basis, just as with respect to any other physical condition.''
\end{quote}
\noindent --- Justice Stewart, for the Court, 417 U.S. at 496--497 n. 20.

\begin{quote}
``The lack of identity between the excluded disability and gender as such under this insurance program becomes clear upon the most cursory analysis. The program divides potential recipients into two groups---pregnant women and nonpregnant persons. While the first group is exclusively female, the second includes members of both sexes.''
\end{quote}
\noindent --- Justice Stewart, for the Court, 417 U.S. at 496--497 n. 20.

\begin{quote}
``There is no evidence in the record that the selection of the risks insured by the program worked to discriminate against any definable group or class in terms of the aggregate risk protection derived by that group or class from the program. There is no risk from which men are protected and women are not. Likewise, there is no risk from which women are protected and men are not.''
\end{quote}
\noindent --- Justice Stewart, for the Court, 417 U.S. at 496--497.

\end{document}
